\chapter{Programmable Hardware I: Architecture and Toolchain}\label{ch:nw:programmable-hardware-i-architecture-and-toolchain}

A \qty{10}{\giga\bit\per\second} Ethernet link delivers a design its bytes at ten billion bits a second; a design that takes them
64 bits at a time must accept a word every \qty{6.4}{\nano\second}, at 156.25~MHz. A design built of registers and logic between them
takes a fixed number of those cycles to decide, and it takes that number every time: no cache miss, no interrupt, no scheduler, no
garbage collector sits between the last byte of a price and the decision. Determinism, rather than raw speed, is what programmable
hardware sells to a trading firm, and it is bought with a different way of writing programs.

This chapter introduces the devices, the languages and the tools, and writes a first design: a field extractor that pulls a price out
of a frame while the frame is still arriving, simulated cycle by cycle in two independent simulators and checked against a Python
model. Chapter~7 builds the trading designs on top of it.

\section{The fabric: logic blocks, memories, transceivers and clocks}

\begin{definition}[Field-programmable gate array, configurable logic block]\label{def:nw:programmable-hardware-i-architecture-and-toolchain:clb}
A \emph{field-programmable gate array}\index{field-programmable gate array} (FPGA) is a chip whose logic and wiring are configured
after manufacture by loading a configuration file: an array of small programmable logic elements, memories, arithmetic blocks and
serial transceivers, joined by programmable routing. A \emph{configurable logic block}\index{configurable logic block} is the array's
repeated logic element: a group of lookup tables, each of which computes any Boolean function of a few inputs by looking it up in a
small memory, with flip-flops to register their outputs and a carry chain for arithmetic.
\end{definition}

\begin{definition}[Block RAM]\label{def:nw:programmable-hardware-i-architecture-and-toolchain:bram}
A \emph{block RAM}\index{block RAM} is a dedicated on-chip memory of a few tens of kilobits, of which an FPGA has hundreds or
thousands, each readable and writable in one clock cycle, typically through two independent ports.
\end{definition}

The resources are concrete. In AMD's UltraScale family each \omterm{def:nw:programmable-hardware-i-architecture-and-toolchain:clb}{configurable logic block} holds eight 6-input lookup tables and sixteen
flip-flops, with an 8-bit carry chain, and each \omterm{def:nw:programmable-hardware-i-architecture-and-toolchain:bram}{block RAM} holds 36 kilobits, usable as one memory or two independent halves. A 6-input
lookup table computes any function of six bits in one level of logic; a function of more bits needs several tables in several levels,
and the depth of that tree of levels, together with the wiring between the tables, sets how fast the design can be clocked. Beside the
logic sit the serial transceivers, which turn the optical module's bit stream into parallel words, and the clocking resources, which
distribute clocks and let parts of the design run at different rates.

\begin{omfigure}
\begin{tikzpicture}[font=\footnotesize,b/.style={draw,rounded corners=2pt,minimum height=0.6cm,align=center,inner sep=3pt}]
\foreach \i in {0,...,5} {\foreach \j in {0,1,2} {\node[b,fill=omDef!12,minimum width=0.9cm] at (1.1*\i,0.9*\j) {CLB};}}
\foreach \j in {0,1,2} {\node[b,fill=omMeth!20,minimum width=0.9cm] at (6.9,0.9*\j) {BRAM};}
\foreach \j in {0,1,2} {\node[b,fill=omThm!15,minimum width=0.9cm] at (8.25,0.9*\j) {DSP};}
\node[b,fill=gray!15,minimum height=2.5cm,minimum width=1.1cm] at (-1.5,0.9) {serial\\ trans-\\ ceivers};
\node[b,fill=gray!8,minimum width=9.6cm] at (3.25,-0.95) {clocking and routing between every block};
\node[b,fill=white,align=left,anchor=west] at (9.2,0.9) {a CLB: 8 lookup\\ tables of 6 inputs,\\ 16 flip-flops,\\ a carry chain};
\end{tikzpicture}

\omcaption{The fabric of an FPGA, schematically: columns of \omterm{def:nw:programmable-hardware-i-architecture-and-toolchain:clb}{configurable logic blocks}, \omterm{def:nw:programmable-hardware-i-architecture-and-toolchain:bram}{block RAMs} and arithmetic (DSP) blocks, serial
transceivers at the edge that turn the optical bit stream into words, and the clocking and routing that join them. Real devices have
hundreds of thousands of lookup tables; the CLB contents are those of AMD's UltraScale family.}\label{fig:nw:programmable-hardware-i-architecture-and-toolchain:fabric}
\end{omfigure}

\section{Describing hardware: register-transfer level and high-level synthesis}

\begin{definition}[Register-transfer level, hardware description language]\label{def:nw:programmable-hardware-i-architecture-and-toolchain:rtl}
A design written at \emph{register-transfer level}\index{register-transfer level} (RTL) describes, for every clock cycle, how the values
held in registers are computed from the registers' previous values and the inputs, through combinational logic. A \emph{hardware
description language}\index{hardware description language} (Verilog, SystemVerilog, VHDL) is a language in which such designs are
written: its statements describe circuits that all operate at once, not instructions executed in order.
\end{definition}

\begin{definition}[Logic synthesis, place and route]\label{def:nw:programmable-hardware-i-architecture-and-toolchain:synth}
\emph{Logic synthesis}\index{logic synthesis} translates an RTL description into a network of the device's primitives (lookup tables,
flip-flops, memories). \emph{Place and route}\index{place and route} assigns each primitive to a location on the chip and chooses the
programmable wires between them; its result is the configuration file loaded into the device.
\end{definition}

\begin{definition}[High-level synthesis]\label{def:nw:programmable-hardware-i-architecture-and-toolchain:hls}
\emph{High-level synthesis}\index{high-level synthesis} (HLS) compiles a function written in C or C++ into an RTL design, choosing its
registers and schedule; the designer steers it with directives such as the pipeline's initiation interval, the number of cycles between
successive inputs (an interval of one accepts a new input every cycle).
\end{definition}

The difference in thinking is the chapter's first lesson. A C++ loop over the eight bytes of a word runs eight times, one after the
other; the same loop in SystemVerilog describes eight copies of its body that all exist and all compute in the same cycle.
\Cref{lst:nw:programmable-hardware-i-architecture-and-toolchain:field} is the combinational heart of the field extractor: for each of the
beat's eight byte lanes it asks whether that byte belongs to the field and, if so, shifts it into the value, all in one cycle. What would
be a data-dependent branch in software becomes a set of multiplexers whose delay does not depend on the data.

\omcode{code/firm/hdlkit/hdl/hdk_field.sv}{19}{29}{The field extractor's combinational logic in SystemVerilog: eight byte lanes examined
in the same cycle; the loop describes eight copies of hardware, not eight iterations.}\label{lst:nw:programmable-hardware-i-architecture-and-toolchain:field}

\section{Timing closure: clock periods, critical paths and pipelining}

\begin{definition}[Critical path, timing closure]\label{def:nw:programmable-hardware-i-architecture-and-toolchain:timing}
The \emph{critical path}\index{critical path} of a synchronous design is its slowest path from one register to the next: the delay of
its logic and wiring plus the registers' own overhead, which sets the shortest clock period the design supports. \emph{Timing
closure}\index{timing closure} is the work of making every register-to-register path fit within the target clock period.
\end{definition}

\begin{definition}[Clock domain crossing]\label{def:nw:programmable-hardware-i-architecture-and-toolchain:cdc}
A \emph{clock domain crossing}\index{clock domain crossing} is a signal passing between parts of a design clocked by unrelated clocks;
it must go through a synchroniser or an asynchronous FIFO, or it will sometimes be sampled while it changes.
\end{definition}

\begin{proposition}[What pipelining buys]\label{prop:nw:programmable-hardware-i-architecture-and-toolchain:pipe}
Let a decision need $L$ levels of logic of delay $t_{\mathrm{lvl}}$ each, and let a register add $t_{\mathrm{reg}}$. Split into $s$
registered stages of at most $\lceil L/s\rceil$ levels, the design's clock period is $T_{\mathrm{clk}} = \lceil L/s\rceil\,t_{\mathrm{lvl}}
+ t_{\mathrm{reg}}$ and its latency is $s\,T_{\mathrm{clk}}$. Pipelining raises the maximum clock $f_{\mathrm{clk}} = 1/T_{\mathrm{clk}}$
and never lowers the latency: $s\,T_{\mathrm{clk}} \ge L\,t_{\mathrm{lvl}} + s\,t_{\mathrm{reg}}$.
\end{proposition}

\begin{proof}
Each stage is one register-to-register path of at most $\lceil L/s\rceil$ levels, so the period is its delay. The decision crosses $s$
stages, one period each, and $s\lceil L/s\rceil \ge L$.
\end{proof}

The proposition explains a design choice that looks backwards: a trading design adds registers, and cycles, to go faster. The line sets
the clock. A \qty{10}{\giga\bit\per\second} link taken 64 bits at a time needs 156.25~MHz; a \qty{25}{\giga\bit\per\second} link at the
same width needs 390.625~MHz. A decision that fits in one cycle at the first clock may need three cycles at the second, and its latency
in nanoseconds then rises a little while the design keeps up with a link two and a half times as fast. \Cref{fig:nw:programmable-hardware-i-architecture-and-toolchain:pipeline}
draws the trade-off with model delays of \qty{0.5}{\nano\second} per level and \qty{0.4}{\nano\second} per register, stated as model values:
a real device's delays come from its timing reports.

\begin{omfigure}
\begin{tikzpicture}
\begin{groupplot}[group style={group size=2 by 1,horizontal sep=1.8cm},width=6.6cm,height=5.0cm,xmin=0.5,xmax=6.5,
  xtick={1,2,3,4,5,6},xlabel={registered stages},tick label style={font=\small},label style={font=\small},grid=major,
  grid style={gray!25},table/col sep=comma]
\nextgroupplot[ymin=0,ymax=800,ylabel={maximum clock (MHz)}]
\addplot[omDef,very thick,mark=*,mark size=1.5pt] table[x=stages,y=fmax_mhz]{figdata/networks/06-programmable-hardware-i-architecture-and-toolchain/pipeline.csv};
\addplot[gray,dashed,domain=0.5:6.5] {156.25};
\addplot[omThm,dashed,domain=0.5:6.5] {390.625};
\node[font=\scriptsize,anchor=south west] at (axis cs:3.6,160) {10G at 64 bits};
\node[font=\scriptsize,anchor=north west,omThm] at (axis cs:4.1,385) {25G at 64 bits};
\nextgroupplot[ymin=0,ymax=10,ylabel={latency at its own clock (ns)}]
\addplot[omMeth,very thick,mark=square*,mark size=1.5pt] table[x=stages,y=latency_ns]{figdata/networks/06-programmable-hardware-i-architecture-and-toolchain/pipeline.csv};
\end{groupplot}
\end{tikzpicture}

\omcaption{Model: a decision of 12 levels of logic split into registered stages (\qty{0.5}{\nano\second} per level,
\qty{0.4}{\nano\second} per register, model values). Left: the maximum clock, against the clocks two \omterm{def:nw:networking-for-trading:ser}{line rates} require; right: the latency
at that clock. Five stages are no better than four, because $\lceil 12/5\rceil = \lceil 12/4\rceil = 3$. Data:
\texttt{fig\_fpga.py}.}\label{fig:nw:programmable-hardware-i-architecture-and-toolchain:pipeline}
\end{omfigure}

\omcode{code/firm/hdlkit/hdl/hdk_cmp.sv}{13}{32}{A comparison in one stage and in two: the two-stage version compares four 16-bit slices,
registers the results, then combines them, one cycle later.}\label{lst:nw:programmable-hardware-i-architecture-and-toolchain:cmp}

\section{Simulation and verification}

\begin{definition}[Testbench, cycle-accurate model]\label{def:nw:programmable-hardware-i-architecture-and-toolchain:tb}
A \emph{testbench}\index{testbench} is the code that drives a design's inputs in simulation and checks its outputs. A \emph{cycle-accurate
model}\index{cycle-accurate model} is a software model of a design that reproduces its outputs cycle by cycle, used as a reference against
which the hardware is checked, and as a fast stand-in for it.
\end{definition}

A design is verified long before it is synthesised, because a synthesised design is slow to build and hard to observe. This book
simulates every design in two independent tools. Verilator compiles the SystemVerilog into a C++ model driven by a C++ \omterm{def:nw:programmable-hardware-i-architecture-and-toolchain:tb}{testbench}: it is
a compiler rather than a traditional simulator, and fast. Icarus Verilog interprets the same source under a SystemVerilog \omterm{def:nw:programmable-hardware-i-architecture-and-toolchain:tb}{testbench}. The
two \omterm{def:nw:programmable-hardware-i-architecture-and-toolchain:tb}{testbenches} read the same stimulus and print the same event lines, and the test compares them cycle for cycle and against a Python
golden model; a disagreement between the simulators is a bug in the design (usually an ambiguity the language permits) before it is a bug
in either tool. The designs are written for the versions installed here (Verilator 4.038, Icarus Verilog 11) and for their newer releases,
and the tests fail, rather than skip, when a simulator is missing.

\begin{table}[ht]
\centering\small
\begin{tabular}{rll}
\toprule
Cycle & Input accepted & Outputs registered at the end of the cycle \\
\midrule
0 & beat 0: bytes \texttt{10}--\texttt{17} & --- \\
1 & beat 1: bytes \texttt{18}--\texttt{1f}, last & --- \\
2 & --- & field \texttt{16171819}; CRC-32 \texttt{f4a7fd67} \\
3 & --- & comparison (one stage): at most the threshold \\
4 & --- & comparison (two stages): at most the threshold \\
\bottomrule
\end{tabular}
\caption{The first frame through \texttt{hdk\_top} in Verilator, no back-pressure: a 16-byte frame of bytes \texttt{10} to \texttt{1f}, the
field at bytes 6 to 9, the threshold equal to the field. The skid buffer registers each accepted beat; the field and the CRC are registered
the cycle after the beat that completes them. Icarus prints the same cycles. Data: \texttt{nw\_fpga.first\_frame}.}\label{tab:nw:programmable-hardware-i-architecture-and-toolchain:cycles}
\end{table}

\begin{omfigure}
\begin{tikzpicture}[font=\footnotesize,b/.style={draw,rounded corners=2pt,minimum height=0.7cm,align=center,inner sep=3pt},>=Stealth]
\node[b,fill=omDef!8] (hls) at (-3.4,0) {C++ with\\ HLS directives};
\node[b,fill=omDef!15] (rtl) at (0,0) {RTL\\ (SystemVerilog)};
\node[b,fill=omMeth!18] (syn) at (2.8,0) {synthesis};
\node[b,fill=omMeth!18] (par) at (5.1,0) {place and\\ route};
\node[b,fill=omMeth!18] (sta) at (7.5,0) {timing\\ analysis};
\node[b,fill=gray!12] (bit) at (10.2,0) {configuration\\ file};
\node[b,fill=omThm!12] (sim) at (4.6,-1.7) {simulation: Verilator, Icarus, against a golden model};
\draw[->] (hls) -- node[above,font=\scriptsize]{HLS} (rtl);
\draw[->] (rtl) -- (syn);
\draw[->] (syn) -- (par);
\draw[->] (par) -- (sta);
\draw[->] (sta) -- node[above,font=\scriptsize]{met} (bit);
\draw[->,omThm] (sta.north) -- ++(0,0.5) -| node[pos=0.25,above,font=\scriptsize]{timing not met: restructure, pipeline} (rtl.north);
\draw[<->,omThm] (rtl.south) |- (sim.west);
\end{tikzpicture}

\omcaption{The toolchain. The design is simulated against its golden model before synthesis; synthesis, placement and routing produce a
design whose timing is analysed; a path that does not meet the clock sends the designer back to the RTL.}\label{fig:nw:programmable-hardware-i-architecture-and-toolchain:flow}
\end{omfigure}

\section{Smart network cards}

\begin{definition}[SmartNIC]\label{def:nw:programmable-hardware-i-architecture-and-toolchain:smartnic}
A \emph{SmartNIC}\index{SmartNIC} is a \omterm{def:nw:networking-for-trading:frame}{network interface card} with its own programmable processing (an FPGA, or processors) between the
network ports and the host, so that packets can be parsed, filtered, timestamped, answered or generated on the card without reaching
the host.
\end{definition}

For a trading firm the FPGA card is where hardware meets the rest of the system: its transceivers terminate the exchange's links, its
logic decodes and decides, and its bus connects it to the host's software, which configures it and receives what it has seen. A design
on such a card reads the frame as it arrives, before the last byte is in, which is what the table of this chapter shows at small scale:
the field is known the cycle after the beat that completes it, while the rest of the frame is still to come. Chapter~7 decides on it.

\section{Tutorial: a field extractor in two simulators}\label{tut:nw:programmable-hardware-i-architecture-and-toolchain:tut}

\begin{tutorial}
\textbf{Goal.} Simulate the first design cycle by cycle in Verilator and Icarus, check both against a Python model under random
back-pressure, and model what pipelining costs. \textbf{End state:} \cref{tab:nw:programmable-hardware-i-architecture-and-toolchain:cycles},
\cref{fig:nw:programmable-hardware-i-architecture-and-toolchain:pipeline} and the green tests of \texttt{firm.hdlkit}.
\begin{tutsteps}
\item \textbf{Read the design.} \texttt{hdl/hdk\_top.sv} chains a skid buffer, the field extractor
  (\cref{lst:nw:programmable-hardware-i-architecture-and-toolchain:field}), a CRC-32 unit and the comparator
  (\cref{lst:nw:programmable-hardware-i-architecture-and-toolchain:cmp}).
\item \textbf{Build both.} \texttt{firm\_hdlkit.build\_verilator} runs \texttt{verilator --cc --exe --build} with the C++ \omterm{def:nw:programmable-hardware-i-architecture-and-toolchain:tb}{testbench};
  \texttt{build\_icarus} runs \texttt{iverilog -g2012} with the SystemVerilog one.
\item \textbf{Run and compare.} \texttt{write\_stim} turns frames into 8-byte beats, \texttt{write\_ready} gives the downstream's
  back-pressure; both simulators print the same event lines, and \texttt{golden} computes the fields, CRCs and comparisons in Python.
\item \textbf{Cycles and clocks.} \texttt{nw\_fpga.first\_frame} gives the table; \texttt{timing} and \texttt{depth\_for} evaluate
  \cref{prop:nw:programmable-hardware-i-architecture-and-toolchain:pipe}, and \texttt{python fig\_fpga.py} writes the chart's data.
\end{tutsteps}
\textbf{What to change next.} Move the field to bytes 60 to 63 and read the cycle it appears in; set the downstream ready to 50\% and count
the cycles the whole stream takes.
\end{tutorial}

\section{Build: the hardware toolkit}\label{bld:nw:programmable-hardware-i-architecture-and-toolchain:hdlkit}

\begin{build}
\bfield{Purpose} The primitives and the verification harness every hardware design of this book uses.
\bfield{Interface} SystemVerilog modules \texttt{hdk\_skid} (skid buffer, parameter \texttt{W}), \texttt{hdk\_field} (\texttt{OFF},
\texttt{LEN}), \texttt{hdk\_crc32}, \texttt{hdk\_cmp} (\texttt{PIPE}) and \texttt{hdk\_top}; \omterm{def:nw:programmable-hardware-i-architecture-and-toolchain:tb}{testbenches} \texttt{tb/hdk\_tb.cpp} (Verilator)
and \texttt{tb/hdk\_tb.sv} (Icarus) printing the same event lines. Python \texttt{firm\_hdlkit}: \texttt{require}, \texttt{beats},
\texttt{write\_stim}, \texttt{write\_ready}, \texttt{build\_verilator}, \texttt{build\_icarus}, \texttt{run\_verilator},
\texttt{run\_icarus}, \texttt{events}, \texttt{golden}, \texttt{CycleModel}.
\bfield{Rules} Beats of 8 bytes, byte 0 in the most significant lane; every output registered; the HDL is accepted by Verilator 4.038 with
\texttt{-Wall} and by Icarus 11 with \texttt{-g2012}, without warnings; a missing simulator raises \texttt{ToolMissing}.
\bfield{Acceptance tests} \texttt{code/firm/hdlkit/tests/}: both simulators identical cycle for cycle and equal to the golden model on random
frames with random back-pressure, for one- and two-stage comparators; every beat accepted once; the latencies in cycles; the beat and golden
functions by hand.
\bfield{Stretch} An asynchronous FIFO for a \omterm{def:nw:programmable-hardware-i-architecture-and-toolchain:cdc}{clock domain crossing} and a test that samples it at unrelated clocks; a wider (512-bit) beat.
\end{build}

\begin{omsources}
\item AMD, \emph{UltraScale Architecture Configurable Logic Block User Guide} (UG574) and \emph{Memory Resources} (UG573); AMD, \emph{Vitis
High-Level Synthesis User Guide} (UG1399).
\item Verilator documentation; Icarus Verilog documentation; IEEE~1800 (SystemVerilog).
\item J.~W. Lockwood and co-authors, ``A Low-Latency Library in FPGA Hardware for High-Frequency Trading (HFT)'', \emph{IEEE Symposium on
High-Performance Interconnects} (2012).
\end{omsources}

\section{Exercises}

\begin{exercise}[$\star$]\label{exo:nw:programmable-hardware-i-architecture-and-toolchain:1}
What clock does a design need to take a \qty{25}{\giga\bit\per\second} link 32 bits at a time, and 128 bits at a time?
\end{exercise}

\begin{exercise}[$\star$]\label{exo:nw:programmable-hardware-i-architecture-and-toolchain:2}
With the model delays, what is the maximum clock of a decision of 12 levels in one stage, and its latency?
\end{exercise}

\begin{exercise}[$\star$]\label{exo:nw:programmable-hardware-i-architecture-and-toolchain:3}
In \cref{tab:nw:programmable-hardware-i-architecture-and-toolchain:cycles}, at which cycle would the field appear if it were at bytes 12 to
15? And at bytes 14 to 17 of a 24-byte frame?
\end{exercise}

\begin{exercise}[$\star\star$]\label{exo:nw:programmable-hardware-i-architecture-and-toolchain:4}
Why does the loop of \cref{lst:nw:programmable-hardware-i-architecture-and-toolchain:field} take one cycle whatever the data, while the same
loop in C++ may not?
\end{exercise}

\begin{exercise}[$\star\star$]\label{exo:nw:programmable-hardware-i-architecture-and-toolchain:5}
In \cref{fig:nw:programmable-hardware-i-architecture-and-toolchain:pipeline}, why is six stages faster than five? What does that say about
counting levels?
\end{exercise}

\begin{exercise}[$\star\star$]\label{exo:nw:programmable-hardware-i-architecture-and-toolchain:6}
What does the skid buffer do that a single register cannot, and what would happen to the stream without it when the downstream stalls?
\end{exercise}

\begin{exercise}[$\star\star\star$]\label{exo:nw:programmable-hardware-i-architecture-and-toolchain:7}
\emph{Coding.} Build \texttt{hdk\_top} with the field at offset 10 and length 8 in both simulators, run the tests' random frames, and check
both against the golden model. How many cycles after the completing beat does the field appear?
\end{exercise}

\begin{exercise}[$\star\star\star$]\label{exo:nw:programmable-hardware-i-architecture-and-toolchain:8}
\emph{Find the flaw.} ``The design passes in Verilator, so it will work on the card; running a second simulator is redundant.''
\end{exercise}

\section{Problem: Eight Cycles}

\begin{problem}[{Weekend problem --- a decision against a line's clock}]\label{pb:nw:programmable-hardware-i-architecture-and-toolchain:1}
A firm's hardware decision (parse a field, compare it with a threshold, check a limit) needs 12 levels of logic. Use the chapter's model:
\qty{0.5}{\nano\second} per level, \qty{0.4}{\nano\second} per register. The design takes the line 64 bits per cycle.

\medskip\noindent\textbf{Part I --- The clocks.}
\begin{enumerate}
\item What clock does a \qty{10}{\giga\bit\per\second} line require at 64 bits a cycle? A \qty{25}{\giga\bit\per\second} line?
\item What is the maximum clock of the decision in one stage?
\item Does it meet either line's clock?
\item How long does a 128-byte frame take to arrive at each line's clock, in cycles and nanoseconds?
\end{enumerate}

\medskip\noindent\textbf{Part II --- Pipelining.}
\begin{enumerate}[resume]
\item How many stages does the \qty{25}{\giga\bit\per\second} line require?
\item What is the maximum clock with that many stages?
\item What is the latency of the decision when the design runs at the line's clock?
\item What is the latency at the design's own maximum clock?
\end{enumerate}

\medskip\noindent\textbf{Part III --- Where the time goes.}
\begin{enumerate}[resume]
\item If the field ends at byte 40 of a 128-byte frame, how many beats must have arrived for the field to be complete?
\item Add the decision at the \qty{25}{\giga\bit\per\second} clock: when is the decision made, counted from the frame's first bit?
\item How much of the frame is still arriving at that moment?
\item Why is it worth deciding before the frame's check sequence arrives, and what is the risk?
\end{enumerate}

\medskip\noindent\textbf{Part IV --- The verdict.}
\begin{enumerate}[resume]
\item State the \emph{named result}: the pipeline depth that meets the \qty{25}{\giga\bit\per\second} clock and the decision's latency, against
  the unpipelined design at its own clock.
\item Would a 128-bit datapath change the answer?
\item What would a real design's timing report replace in this model?
\item Why does pipelining never reduce latency?
\item Why may the firm accept a longer decision latency for a faster line?
\item How would you verify the pipelined design against the unpipelined one?
\item What does a \omterm{def:nw:programmable-hardware-i-architecture-and-toolchain:cdc}{clock domain crossing} add if the decision runs at a clock other than the line's?
\item In one sentence: what does the line's rate decide about the design?
\end{enumerate}
\end{problem}

\section{Interview questions}

\begin{interviewq}[$\star$ \iqroles{developer}]\label{iq:nw:programmable-hardware-i-architecture-and-toolchain:1}
What is an FPGA, and why do trading firms use them?
\end{interviewq}

\begin{interviewq}[$\star\star$ \iqroles{developer}]\label{iq:nw:programmable-hardware-i-architecture-and-toolchain:2}
What is \omterm{def:nw:programmable-hardware-i-architecture-and-toolchain:timing}{timing closure}? What do you do when a path fails timing?
\end{interviewq}

\begin{interviewq}[$\star\star$ \iqroles{developer}]\label{iq:nw:programmable-hardware-i-architecture-and-toolchain:3}
Pipelining adds cycles. How can it make a design faster?
\end{interviewq}

\begin{interviewq}[$\star\star$ \iqroles{developer}]\label{iq:nw:programmable-hardware-i-architecture-and-toolchain:4}
RTL or \omterm{def:nw:programmable-hardware-i-architecture-and-toolchain:hls}{high-level synthesis} for a trading design? Argue both sides.
\end{interviewq}

\begin{interviewq}[$\star\star\star$ \iqroles{developer}]\label{iq:nw:programmable-hardware-i-architecture-and-toolchain:5}
How would you verify a hardware feed parser before it runs against a live exchange?
\end{interviewq}

\begin{interviewq}[$\star\star$ \iqroles{developer}]\label{iq:nw:programmable-hardware-i-architecture-and-toolchain:6}
What is a \omterm{def:nw:programmable-hardware-i-architecture-and-toolchain:cdc}{clock domain crossing}, and what goes wrong if it is not handled?
\end{interviewq}
